\section{Discussion, Limitations and Future Work}
\label{sec:discussion}

The experiments show that completion depends on both course structure and environmental conditions: the two reference controllers clear the short clean circuit but separate on the longer track and under current, and success on clean tracks does not imply robustness to disturbances. The coordination study shows that a fully distributed leader-follower layer, using only legal onboard estimates and an acoustic-inspired channel, removes the physical interaction that an uncoordinated heterogeneous convoy suffers, with the one-gate margin doing so at the lowest time cost. Throughout, controller-local progression and independent referee scoring are kept separate, and their occasional disagreement is logged rather than reconciled, which is what lets the same contract score arbitrary controllers.

\subsection{Limitations}
\label{sec:limitations}

The evaluation is deliberately scoped, and several boundaries should be read with the results. The quantitative current sweep is limited to Horseshoe Bay; the medium and strong profiles are available on all three circuits but only Horseshoe Bay is measured. The coordination validation is limited to one clean track, three vehicles, two start gaps and three seeds. Both reference controllers are rule-based and interpretable rather than optimized, and no learning-based controller is reported. All results are in simulation: no real underwater-vehicle experiment is included, and the numbers depend on the HoloOcean vehicle and sensor models. The inter-vehicle acoustic channel is exercised only as the substrate for coordination; its packet-loss, range, latency and freshness parameters are not independently characterized. The strong-current condition acts primarily as a saturation regime that limits forward progress rather than a controllable disturbance. Finally, the two-parameter beacon-noise model preserves the numerical behavior of the evaluated single-scalar model exactly (verified by a byte-level equivalence test) rather than introducing newly calibrated physical sensor statistics. None of these boundaries invalidates the reported findings; they mark where the released benchmark invites extension, which the configuration-driven design is built to support.

\subsection{Future Work}

We see four concrete next steps.

\emph{(i) Cooperative fleet strategies beyond leader-follower convoying.} The leader-follower coordination of Section~\ref{sec:coordination} is validated on clean Horseshoe Bay (Section~\ref{sec:eval_holoocean_coord}). A natural extension covers other tracks, currents, obstacles and larger fleets, and then formation keeping, cooperative overtaking, dynamic role assignment and robustness to uncertain relative localization or heavier acoustic loss. Such strategies must respect the low bandwidth and high latency of acoustic communication~\cite{stojanovic2009underwater} described in the AUV-formation literature~\cite{yang2021survey,yan2023formation}.

\emph{(ii) Competitive multi-team and heterogeneous-vehicle evaluation.} Extend the cooperative fleet setting to multiple teams or vehicle families with different dynamics and sensing suites, competing under the same referee and observation contract. This requires team-isolated scoring, ranking across teams, controlled interaction policies and normalization rules that compare heterogeneous platforms without hiding their physical differences.

\emph{(iii) Current compensation and broader disturbance validation.} The three official circuits already expose the graded \texttt{medium} and \texttt{strong} profiles of \eqref{eq:current_sum} to \eqref{eq:vortex}; the present quantitative sweep is limited to Horseshoe Bay. Controllers can be evaluated on the remaining circuits and can explicitly reject these disturbances. One legal direction is a cascaded design with an inner loop that servos DVL-measured body velocity to the guidance setpoint, beneath outer beacon and vision guidance, without access to environmental-current telemetry.

\emph{(iv) Learning-based controllers.} The official observation contract can serve directly as a reinforcement-learning interface. The state contains received beacon range, bearing and elevation; depth; IMU; body-frame DVL; and camera data. The normalized command of \eqref{eq:command} is the action. A reward can combine locally estimated sequential progress with the time and collision penalties of \eqref{eq:penalized}, while official progress remains evaluator-only. Training and evaluation should retain the HoloOcean vehicle and contact dynamics used for the reported benchmark results. This direction follows work on champion-level drone racing against a simulated opponent~\cite{kaufmann2023swift} and robot-learning environments~\cite{makoviychuk2021isaacgym,brockman2016openaigym} without weakening the official observation contract.
